\begin{table}[htbp]
\centering
\caption{Composition of employment by educational attainment}
\label{tab:latam_country_education_shares}
\small
\resizebox{\textwidth}{!}{%
\begin{tabular}{lrrrrrr}
\toprule
Economy & Low initial & Low final & Middle initial & Middle final & High initial & High final \\
\midrule
Argentina & 36.0 & 29.9 & 41.0 & 44.6 & 22.9 & 25.5 \\
Bolivia & 57.4 & 47.4 & 36.5 & 44.9 & 6.1 & 7.7 \\
Brazil & 42.6 & 33.0 & 40.5 & 45.3 & 16.9 & 21.6 \\
Chile & 31.9 & 19.8 & 53.9 & 55.3 & 14.1 & 24.8 \\
Colombia & 46.1 & 34.9 & 40.2 & 35.6 & 13.7 & 29.6 \\
Costa Rica & 54.5 & 44.9 & 32.3 & 38.6 & 13.1 & 16.5 \\
Dominican Republic & 55.3 & 48.4 & 31.4 & 36.6 & 13.4 & 15.0 \\
Ecuador & 44.4 & 48.4 & 38.1 & 36.6 & 17.5 & 15.0 \\
El Salvador & 75.2 & 72.6 & 20.6 & 21.0 & 4.2 & 6.4 \\
Guatemala & 72.1 & 79.1 & 23.1 & 15.0 & 4.7 & 5.9 \\
Mexico & 61.7 & 53.3 & 21.6 & 25.5 & 16.7 & 21.2 \\
Paraguay & 56.1 & 46.7 & 31.3 & 36.0 & 12.6 & 17.3 \\
Peru (Lima and Callao) & 17.5 & 16.0 & 51.4 & 49.4 & 31.2 & 34.6 \\
Uruguay & 64.1 & 60.3 & 22.3 & 22.9 & 13.6 & 16.8 \\
\bottomrule
\end{tabular}
}
\begin{minipage}{0.98\textwidth}
\footnotesize
\textit{Note:} Percent of employed workers in each education group. Initial and final refer to the years reported in Table~\ref{tab:latam_country_income_descriptives}.
\end{minipage}
\end{table}
